% ==================================================================
% Conscious Point Physics:
% The Nexus as Superdeterministic Mechanism:
% Bell Correlations, Hidden Variables, and the 600-Cell Angular Signature
% CPP Companion SD#1 — Version 1
%
% Position paper / open problem statement
% Connects CPP to Hossenfelder (2024) superdeterminism programme
% GitHub: CPP/companions/cpp_sd1_nexus_superdeterminism_v1.tex
% ==================================================================

\documentclass[11pt,a4paper]{article}
\usepackage[margin=1in]{geometry}
\usepackage[utf8]{inputenc}
\usepackage[T1]{fontenc}
\usepackage{amsmath,amssymb,amsthm}
\usepackage{graphicx}
\usepackage{caption,float,booktabs,parskip,xcolor}
\usepackage[numbers,sort&compress]{natbib}
\usepackage{hyperref}
\hypersetup{colorlinks=true,linkcolor=blue,urlcolor=cyan,citecolor=red}

\newtheorem{theorem}{Theorem}[section]
\newtheorem{proposition}[theorem]{Proposition}
\newtheorem{remark}[theorem]{Remark}
\newtheorem{openprob}[theorem]{Open Problem}
\newtheorem{prediction}[theorem]{Prediction}

\newcommand{\lP}{l_P}
\newcommand{\tP}{t_P}
\newcommand{\SSV}{\mathrm{SSV}}
\newcommand{\hDP}{\mathrm{hDP}}
\newcommand{\qCP}{\mathrm{qCP}}
\newcommand{\eCP}{\mathrm{eCP}}
\newcommand{\phig}{\varphi}
\newcommand{\eps}{\varepsilon}

\title{\textbf{Conscious Point Physics:\\
The Nexus as Superdeterministic Mechanism:\\
Bell Correlations, Hidden Variables,\\
and the 600-Cell Angular Signature}\\[6pt]
{\large Companion SD\#1 --- Version 1}\\[4pt]
{\normalsize \textit{Position Paper and Open Problem Statement}}}

\author{Thomas Lee Abshier, ND\\
Hyperphysics Institute\\
\url{https://hyperphysics.com}\\
\texttt{drthomas007@protonmail.com}}

\date{23 March 2026}

\begin{document}
\maketitle

\begin{abstract}
Hossenfelder~\citep{hossenfelder2024} argues that quantum mechanics is
not fundamental but is the statistical average of an underlying
deterministic hidden-variables theory (superdeterminism), and that
the apparent impossibility of this programme is a circular argument
maintained by groupthink rather than by physical necessity.

Conscious Point Physics (CPP) provides a concrete instantiation of
precisely this programme.  The DP~Sea configurations are the hidden
variables; the 600-cell lattice is the deterministic substrate; and
the Nexus --- an atemporal global constraint on DI-bit conservation ---
is the mechanism that produces the measurement-setting correlations
required by superdeterminism.

This companion paper does four things.
\textbf{(1)}~It states the open problem precisely in CPP language: the
Nexus must correlate the apparatus orientation $\theta$ with the
hidden-variable DP~Sea state $\lambda$ at the time of entanglement,
and the required correlation function has not yet been derived from
first principles.
\textbf{(2)}~It establishes the connection to Hossenfelder's
framework in detail.
\textbf{(3)}~It derives order-of-magnitude estimates for the magnitude
of the deviation from standard quantum mechanics: the dominant term
is the Nexus participation fraction $\eps \approx N_{\rm part} /
N_{\rm app} \sim 10^{-26}$ for a two-particle Bell test against a
macroscopic detector, and the correction to the CHSH parameter is
$\delta S \approx \eps \cdot 2\sqrt{2} \sim 10^{-25}$.
\textbf{(4)}~It identifies the key experimental signature: the
correction has a specific angular dependence set by the 600-cell
icosahedral symmetry, with predicted structure at angles
$\{31.7°, 36°, 45°, 58.3°, 67.5°, 72°, 120°\}$ --- the
golden-ratio angles of the $H_4$ polytope.  This angular signature
is falsifiable even though the magnitude is far below current
experimental precision.
\end{abstract}

\tableofcontents
\newpage

%=====================================================================
\section{Introduction: Why This Open Problem Matters}
%=====================================================================

The CPP programme has now derived the quantum mechanical, electroweak,
and strong sectors of the Standard Model from the 600-cell lattice
\citep{cpp_qm_series,cpp_ew_unified,cpp_ss_unified}.  The QM series
derived the Schrödinger equation, Born rule, Bell inequality
violations, and the Lindblad master equation from DP~Sea dynamics.
In each case, quantum mechanics is an \emph{emergent} statistical
description --- not a fundamental law.

This raises an immediate question that the QM series deferred: if QM
is statistical, what is the underlying deterministic structure, and
how does CPP avoid Bell's theorem?

Bell's theorem proves that no \emph{locally realistic} hidden-variable
theory can reproduce all predictions of quantum mechanics, given the
assumption of \emph{statistical independence} --- that the measurement
settings $\theta_A, \theta_B$ are independent of the hidden variables
$\lambda$.  Superdeterminism is the logical complement: if statistical
independence is violated, Bell's theorem does not apply, and a local
hidden-variable theory becomes possible.

CPP is a superdeterministic theory.  The hidden variables are the
DP~Sea configurations; the Nexus is the mechanism that correlates
them with apparatus orientations.  This paper states precisely what
CPP must derive to make this programme complete, and what the theory
predicts in the meantime.

\subsection{Connection to Hossenfelder (2024)}

Hossenfelder \citep{hossenfelder2024} makes three central arguments:

\begin{enumerate}
\item A preferred time direction (preferred slicing of spacetime)
  eliminates causality paradoxes from faster-than-light signaling.
\item Quantum mechanics is a statistical theory over hidden variables
  (superdeterminism), and its resemblance to the Liouville equation
  of statistical mechanics is evidence for this.
\item If hidden variables are real, deviations from standard QM
  allow FTL correlations.  Quantum computers will produce anomalous
  data that AI will identify as consistent with a hidden-variable
  theory.
\end{enumerate}

CPP has specific counterparts to each of these.  The DP~Sea rest
frame provides the preferred slicing.  The DP~Sea phase-space
distribution plays the role of the Liouville distribution.  The
Nexus provides the superdeterministic correlations.  This is not
an analogy; CPP provides the geometric structure that Hossenfelder's
programme identifies as necessary but does not specify.

%=====================================================================
\section{The CPP Hidden-Variable Structure}
%=====================================================================

\subsection{Hidden variables in CPP}

In CPP, every measurement outcome is determined by the configuration
of the DP~Sea at the time of the interaction.  The hidden variable
$\lambda$ for an entangled pair is the joint DP~Sea state of the
two particles and all the hDP structures mediating their interaction:
\begin{equation}
\lambda = \{\rho_i, \phi_i, s_i\}_{i \in \mathcal{N}_{\rm pair}},
\label{eq:hv}
\end{equation}
where $\rho_i$ is the DI-bit density, $\phi_i$ is the phase, and
$s_i$ is the spin orientation at each Grid~Point $i$ in the causal
neighbourhood $\mathcal{N}_{\rm pair}$ of the pair.  Standard QM
averages over this distribution.

\subsection{Why CPP is local}

The SSV mechanism is local: the SSV at each Grid~Point depends only
on the DP~Sea density and gradient at that point and its 12 nearest
neighbours (the 600-cell vertex coordination number).  No
instantaneous action at a distance is required for any individual
interaction.

The apparent non-locality of entanglement arises from the fact that
the hidden variables $\lambda$ were set at the time of entanglement
to be correlated across both particles.  The Nexus enforces this
correlation atemporally, but the actual causal chain is local: the
two particles shared a common causal past in which their DP~Sea
states were jointly constrained.

\subsection{The statistical independence assumption in CPP}

Standard Bell analysis assumes:
\begin{equation}
\rho(\lambda | \theta_A, \theta_B) = \rho(\lambda),
\label{eq:stat_indep}
\end{equation}
i.e., the hidden-variable distribution does not depend on the
measurement settings.  In CPP, this assumption fails because:

\begin{enumerate}
\item The apparatus is a macroscopic collection of $\sim 10^{26}$ eCPs.
\item The apparatus orientation $\theta$ is a collective degree of
  freedom of these eCPs.
\item The Nexus enforces global DI-bit conservation:
  $\sum_{i} \delta b_i = 0$ across the entire lattice at all times.
\item Therefore the DP~Sea state $\lambda$ of the particle pair is
  not independent of the DP~Sea state of the apparatus, because both
  are constrained by the same global Nexus condition.
\end{enumerate}

This is the precise statement of why CPP is superdeterministic.  It
does not require any spooky action; it requires only that the entire
lattice is globally constrained by the Nexus, and that the
measurement apparatus is part of the same lattice as the particles.

%=====================================================================
\section{The Open Problem: Deriving the Nexus Correlation Function}
\label{sec:openprob}
%=====================================================================

\subsection{What must be derived}

Let $\theta_A$ be the orientation of Alice's detector and $\theta_B$
be the orientation of Bob's detector.  The Nexus correlation function
$K$ is defined by:
\begin{equation}
\rho(\lambda | \theta_A, \theta_B)
= \rho(\lambda)\left[1 + K(\lambda, \theta_A, \theta_B)\right],
\label{eq:nexus_corr}
\end{equation}
where $K \to 0$ in the classical limit (macroscopic apparatus,
$\hbar \to 0$).

The open problem is to derive $K$ from CPP first principles, i.e.,
from the Nexus constraint on the 600-cell lattice.  Specifically:
\begin{itemize}
\item What is the functional form of $K(\lambda, \theta_A, \theta_B)$?
\item What is the magnitude of $K$ as a function of the apparatus
  size $N_{\rm app}$ and particle count $N_{\rm part}$?
\item What is the angular dependence of $K$ on $\theta_A - \theta_B$?
\item At what experimental scale would $K$ become detectable?
\end{itemize}

\begin{openprob}[Nexus correlation function]
\label{op:nexus}
Derive the correlation function $K(\lambda, \theta_A, \theta_B)$
in equation~\eqref{eq:nexus_corr} from the Nexus constraint
$\sum_i \delta b_i = 0$ on the 600-cell lattice.  Show that:
(a)~$K$ vanishes as $N_{\rm app} \to \infty$ at fixed $N_{\rm part}$
(macroscopic limit recovers standard QM); (b)~$K$ has an angular
dependence set by the $H_4$ symmetry group of the 600-cell;
(c)~$K$ reduces to the CPP Born rule (QM series, companion C3
\citep{cpp2040c}) when averaged over the apparatus degrees of freedom.
\end{openprob}

\subsection{What is known: constraints on K}

Although the full derivation is open, several constraints on $K$
follow from the CPP structure already established:

\begin{proposition}[Constraints on the Nexus correlation function]
\label{prop:constraints}
The function $K(\lambda, \theta_A, \theta_B)$ satisfies:
\begin{align}
\int K(\lambda, \theta_A, \theta_B)\,\rho(\lambda)\,d\lambda &= 0
\label{eq:K_zero_mean} \\
|K(\lambda, \theta_A, \theta_B)| &\lesssim \frac{N_{\rm part}}{N_{\rm app}}
\label{eq:K_magnitude} \\
K(\lambda, \theta_A, \theta_B) &= K_0(\lambda) \cdot
f_{H_4}(\theta_A - \theta_B)
\label{eq:K_factored}
\end{align}
where $K_0(\lambda)$ is a function of the hidden variables alone and
$f_{H_4}$ is a function with the icosahedral $H_4$ symmetry of the
600-cell.
\end{proposition}

\begin{proof}[Argument for~\eqref{eq:K_zero_mean}]
Averaging over hidden variables must recover the standard QM
distribution: $\int \rho(\lambda|\theta)\,d\lambda = 1$.\qed
\end{proof}

\begin{proof}[Argument for~\eqref{eq:K_magnitude}]
The Nexus distributes the DI-bit conservation constraint over all
$N_{\rm total}$ Grid~Points in the lattice.  The pair contributes
$N_{\rm part}$ and the apparatus contributes $N_{\rm app}$.  The
fractional influence of the pair on the apparatus state, and vice
versa, scales as $N_{\rm part}/N_{\rm total} \approx
N_{\rm part}/N_{\rm app}$ when $N_{\rm app} \gg N_{\rm part}$.\qed
\end{proof}

\begin{proof}[Argument for~\eqref{eq:K_factored}]
The 600-cell's adjacency matrix has $H_4$ symmetry (order 14400).
The Nexus constraint is covariant under $H_4$ rotations.  The
apparatus orientation $\theta$ enters only through the collective
variable $\theta_A - \theta_B$, which the Nexus ``sees'' through
the relative orientation of the two detector subgraphs in the
lattice.  Therefore the angular part of $K$ must be an $H_4$-covariant
function of $\theta_A - \theta_B$.\qed
\end{proof}

%=====================================================================
\section{Order-of-Magnitude Estimates}
\label{sec:magnitude}
%=====================================================================

\subsection{The Nexus participation fraction}

For a standard loophole-free Bell test with two photons measured by
macroscopic polarisers:
\begin{equation}
N_{\rm part} = 2 \quad (\text{two photons}), \qquad
N_{\rm app} \sim 10^{26} \quad (\text{atoms in detector}).
\label{eq:N_values}
\end{equation}
The Nexus participation fraction is:
\begin{equation}
\eps_{\rm Nexus} = \frac{N_{\rm part}}{N_{\rm app}}
\approx \frac{2}{10^{26}} = 2 \times 10^{-26}.
\label{eq:eps}
\end{equation}

The correction to the CHSH parameter $S$ (standard QM maximum:
$S_{\rm QM} = 2\sqrt{2} \approx 2.828$) is:
\begin{equation}
|\delta S|_{\rm CPP} \approx \eps_{\rm Nexus} \cdot S_{\rm QM}
\approx 5.7 \times 10^{-26}.
\label{eq:deltaS}
\end{equation}

This is 22 orders of magnitude below the current best experimental
precision of $|\delta S|_{\rm exp} \sim 10^{-3}$ \citep{li2018}.
The deviation is undetectable in current Bell tests but is
\emph{not zero} --- it is a finite, calculable prediction of CPP.

\subsection{Comparison with the naive Planck-scale estimate}

A simpler estimate uses the ratio of the Planck length to the de
Broglie wavelength of the particles:
\begin{equation}
\eps_{\rm Planck} \sim \left(\frac{\lP}{\lambda_{\rm dB}}\right)^2.
\label{eq:eps_planck}
\end{equation}
For optical photons ($\lambda_{\rm dB} = 800$~nm):
$\eps_{\rm Planck} \sim 4 \times 10^{-58}$.
This is 32 orders smaller than the participation-fraction estimate.

The participation-fraction estimate~\eqref{eq:eps} is more
conservative and more physically motivated: it counts how many
lattice sites are constrained by the particle vs.\ the apparatus,
rather than assuming the correction scales with wavelength.  The
two estimates bound the true answer from below ($\eps_{\rm Planck}$)
and above ($\eps_{\rm Nexus}$); the full derivation
(Open Problem~\ref{op:nexus}) will determine which is correct.

\subsection{Quantum processor advantage}

Quantum computers offer a different regime.  In a quantum processor
with $N_q$ qubits, all $N_q$ qubits form a single quantum system
simultaneously participating in the Nexus constraint.  The effective
participation fraction is enhanced:
\begin{equation}
\eps_{\rm QP} = \frac{N_q}{N_{\rm substrate}},
\label{eq:eps_QP}
\end{equation}
where $N_{\rm substrate}$ is the number of atoms in the dilution
refrigerator substrate ($\sim 10^{26}$).  For $N_q = 1000$ qubits:
$\eps_{\rm QP} \sim 10^{-23}$.  For $N_q = 10^6$: $\eps_{\rm QP}
\sim 10^{-20}$.

More importantly, Bell-inequality tests on quantum processors can
scan $\theta_A - \theta_B$ continuously with high resolution.
The CPP prediction is not just a scalar shift in $S$ but a specific
angular dependence, described in the next section.

%=====================================================================
\section{The 600-Cell Angular Signature}
\label{sec:angular}
%=====================================================================

\subsection{$H_4$ symmetry and special angles}

The 600-cell has symmetry group $H_4$ (order 14400), the highest-order
hyperoctahedral group.  It contains the following subgroups relevant
to the angular structure of $f_{H_4}(\theta)$:

\begin{table}[H]
\centering
\caption{600-cell special angles and their geometric origin.
  These are predicted to be the locations of extrema in the CPP
  correction function $f_{H_4}(\theta_A - \theta_B)$.}
\label{tab:angles}
\begin{tabular}{@{}lcc@{}}
\toprule
Origin & Angle & Geometric structure \\
\midrule
$\arctan(1/\phig)$ & $31.72°$ & Golden-ratio base angle \\
Pentagon inscribed & $36.00°$ & $\pi/5$ ($H_2$ subgroup) \\
CHSH optimal & $45.00°$ & $\pi/4$ (standard QM maximum) \\
$\arctan(\phig)$ & $58.28°$ & Golden-ratio complement \\
$\arccos(1/\phig^2)$ & $67.54°$ & 600-cell radial angle \\
Pentagonal symmetry & $72.00°$ & $2\pi/5$ ($H_2$ subgroup) \\
Triangular symmetry & $120.00°$ & $2\pi/3$ (SU(2) generator spacing) \\
\bottomrule
\end{tabular}
\end{table}

\subsection{Predicted form of the angular correction}

The leading-order CPP correction to the CHSH correlation function
$E(\theta_A, \theta_B) = -\cos(\theta_A - \theta_B)$ takes the form:
\begin{equation}
E_{\rm CPP}(\theta_A, \theta_B)
= -\cos(\theta_A - \theta_B)
  + \eps \cdot f_{H_4}(\theta_A - \theta_B) + O(\eps^2),
\label{eq:E_CPP}
\end{equation}
where $\eps \sim \eps_{\rm Nexus}$ from~\eqref{eq:eps} and
$f_{H_4}$ is an $H_4$-covariant periodic function of
$\theta = \theta_A - \theta_B$.

At leading order, $f_{H_4}$ is the lowest-order $H_4$-invariant
harmonic on the circle.  The $H_4$ group projects onto the plane
with 5-fold symmetry as its leading component (from the icosahedral
$H_3$ subgroup of $H_4$).  The minimal such function is:
\begin{equation}
f_{H_4}(\theta) \approx A_5 \cos(5\theta) + A_{10}\cos(10\theta)
+ A_3 \cos(3\theta) + \ldots,
\label{eq:f_H4}
\end{equation}
where $A_5, A_{10}, A_3$ are amplitudes to be derived from the Nexus
correlation (Open Problem~\ref{op:nexus}).

The 5-fold component $\cos(5\theta)$ has extrema at $\theta = 0°,
36°, 72°, 108°, 144°$ --- the pentagon angles.  The 3-fold component
$\cos(3\theta)$ has extrema at $\theta = 0°, 60°, 120°$ --- the
triangle angles.  Their combination produces structure at all the
golden-ratio angles of Table~\ref{tab:angles}.

\begin{prediction}[600-cell angular signature in Bell tests]
\label{pred:angular}
A precision scan of the CHSH correlation function $E(\theta)$ over
$\theta \in [0°, 180°]$ will reveal a residual after subtracting
$-\cos\theta$ that has:
\begin{enumerate}
\item Extrema at the golden-ratio angles of Table~\ref{tab:angles},
  particularly $\theta = 36°$, $72°$, and the composite at $31.7°$
  and $58.3°$.
\item Amplitude $|\eps \cdot f_{H_4}| \leq \eps_{\rm Nexus}
  \approx 2 \times 10^{-26}$ for standard Bell tests,
  scaling as $N_{\rm part}/N_{\rm app}$.
\item No deviation at $\theta = 45°$ (CHSH optimal angle): the
  correction vanishes there because $45°$ is not a special angle
  of the 600-cell, so $f_{H_4}(45°) \approx 0$.
\end{enumerate}
This signature is \emph{qualitatively} different from standard QM,
which predicts a pure $-\cos\theta$ dependence, and from other
hidden-variable theories, which predict deviations without the
specific icosahedral structure.
\end{prediction}

\begin{remark}[Falsifiability vs.\ detectability]
The amplitude $\sim 10^{-26}$ is far below current experimental
sensitivity.  However, falsifiability does not require immediate
detectability.  The prediction is falsifiable because:
(a)~if CPP is correct, the angular structure \emph{must} appear as
precision improves; (b)~if it does not appear when precision reaches
$\sim 10^{-26}$, CPP is ruled out; (c)~if it does appear, the
specific icosahedral pattern distinguishes CPP from all other
hidden-variable theories.

This is the same standard applied to gravitational wave detection
(decades of effort) and the Higgs boson (decades of accelerator
progress).  The relevant question is whether the prediction is
\emph{specific} and \emph{non-trivial} --- it is both.
\end{remark}

%=====================================================================
\section{The Measurement Apparatus as a CPP Structure}
\label{sec:apparatus}
%=====================================================================

A complete derivation of $K$ requires modelling the measurement
apparatus as a CPP structure.  This is the most technically demanding
part of the programme.

\subsection{What the apparatus is in CPP}

A polarisation analyser or spin detector is a macroscopic arrangement
of $\sim 10^{26}$ eCPs.  Its orientation $\theta$ is a collective
degree of freedom: the macroscopic alignment of the crystalline or
magnetic structure sets a preferred direction in the DP~Sea.

In the 600-cell lattice, a macroscopic orientation $\theta$ corresponds
to a preferred axis in 3D space, which lifts the $H_4$ symmetry of
the vacuum DP~Sea to a subgroup.  The specific subgroup depends on
$\theta$: at $\theta = 72°$ the full $H_2$ (pentagonal) symmetry is
preserved; at generic angles a lower symmetry remains.

\subsection{Decoherence and the classical limit}

The CPP hidden variables decohere on the timescale $\tau_{\rm dec}
\sim \hbar / (N_{\rm app} k_B T)$, where $T$ is the apparatus
temperature.  For a room-temperature detector: $\tau_{\rm dec} \sim
10^{-13}$~s, much shorter than the measurement time.  This is why
the apparatus acts classically: its quantum DP~Sea fluctuations
average to zero on the measurement timescale.

The quantum system (the entangled pair) has decoherence time
$\tau_{\rm dec}^{\rm pair} \gg \tau_{\rm meas}^{\rm pair}$, which is
why quantum effects persist.

\subsection{The apparatus-particle correlation: what is needed}

The full derivation of $K$ requires:
\begin{enumerate}
\item A CPP treatment of the macroscopic apparatus as a coherent
  sum of $N_{\rm app}$ eCPs with collective orientation $\theta$.
\item A calculation of how the Nexus constraint couples this
  collective variable to the hidden variable $\lambda$ of the pair.
\item An integration over apparatus degrees of freedom to obtain
  $\rho(\lambda | \theta_A, \theta_B)$.
\end{enumerate}

Steps 1 and 2 are new calculations not yet in the CPP programme.
Step 3 is the standard statistical mechanics calculation once the
CPP expression for step 2 is in hand.

%=====================================================================
\section{Relation to the Hossenfelder Programme}
\label{sec:hossenfelder}
%=====================================================================

\subsection{Three alignments}

Hossenfelder identifies three claims:

\textbf{Claim~1 (preferred slicing).}
CPP has the DP~Sea rest frame as its preferred slicing.
Lorentz symmetry is emergent (SR companion C2 \citep{abshier_sr}).
The preferred frame is real but its effects are suppressed by the
ratio of laboratory scales to cosmological scales, making it
undetectable at current precision --- consistent with the observed
Lorentz invariance of all current experiments.

\textbf{Claim~2 (QM as statistical mechanics).}
CPP derives the Schrödinger equation (QM series, companion 2040a),
the Born rule (companion C3), and the Lindblad master equation
(companion 2040d) from DP~Sea statistics.  The resemblance of
the Dirac--von Neumann equation to the Liouville equation is
exact in CPP: both are continuity equations in the same DP~Sea
phase space, at different levels of coarse-graining.

\textbf{Claim~3 (deviations from QM enable FTL correlations).}
In CPP, the Nexus operates atemporally across the entire lattice.
Deviations from standard QM arising from the Nexus correlation
function $K$ allow measurement outcomes on one side to correlate
with outcomes on the other side through the pre-established
DP~Sea state --- faster than any signal, but not in violation of
causality because the correlations were set at entanglement time
in the DP~Sea rest frame (Claim~1).

\subsection{What CPP adds to Hossenfelder's programme}

Hossenfelder's programme is a philosophical framework identifying
what a complete theory must accomplish.  CPP provides:
\begin{itemize}
\item A specific geometric structure: the 600-cell lattice with its
  six adjacency eigenvalues.
\item A specific hidden-variable space: the DP~Sea phase-space
  distribution over Grid~Point density $\rho_i$ and phase $\phi_i$.
\item A specific superdeterministic mechanism: the Nexus global
  DI-bit conservation.
\item A specific symmetry for the angular correction: $H_4$
  (icosahedral), not arbitrary.
\item A specific prediction: the angular signature of
  Prediction~\ref{pred:angular}.
\end{itemize}

The key contribution is that CPP converts Hossenfelder's programme
from a no-go theorem critique into a concrete theory with specific
predictions.

\subsection{Hossenfelder's AI prediction and CPP}

Hossenfelder predicts that AI analysing quantum computer data will
identify hidden-variable signatures ``in the coming decade.''  CPP
specifies what those signatures should look like:

\begin{enumerate}
\item Deviations from the Born rule statistics in high-gate-count
  quantum circuits, scaling as $\eps_{\rm QP} \sim N_q / N_{\rm sub}$.
\item Anomalous correlations in multi-qubit Bell tests with angular
  dependence fitting $f_{H_4}(\theta)$ rather than pure $\cos\theta$.
\item The anomalies grow with $N_q$ (more qubits = larger participation
  fraction), distinguishing CPP from noise or systematic error.
\item The angular structure has 5-fold symmetry ($\cos 5\theta$) as
  the leading term, distinct from any other known hidden-variable theory.
\end{enumerate}

If quantum computers show anomalies with these specific properties,
CPP is confirmed.  If they show anomalies without the icosahedral
structure, CPP is ruled out in favour of a different hidden-variable
theory.  If they show no anomalies to precision $\sim 10^{-26}$, CPP
is ruled out entirely.

%=====================================================================
\section{Consolidated Open Problems}
%=====================================================================

\begin{openprob}[Derive the Nexus correlation function $K$]
\label{op:K}
(Restating Open Problem~\ref{op:nexus} with formal precision.)
Derive
\[
K(\lambda, \theta_A, \theta_B) = \frac{\rho(\lambda|\theta_A,\theta_B)
- \rho(\lambda)}{\rho(\lambda)}
\]
from the CPP Nexus constraint $\sum_i \delta b_i = 0$ on the
600-cell lattice.  Establish the three properties of
Proposition~\ref{prop:constraints} rigorously.
This is the primary open problem of the CPP quantum-foundations
programme; it is to the QM sector what deriving
$\eta_W \approx 1.57\times10^{-17}$ (Open Problem EW-1) is to the
electroweak sector.
\end{openprob}

\begin{openprob}[Apparatus model in CPP]
\label{op:apparatus}
Construct a CPP model of a macroscopic measurement apparatus:
a polariser/spin-detector as $\sim10^{26}$ eCPs with collective
orientation $\theta$.  Derive the effective DP~Sea density
$\bar\rho_{\rm app}(\theta)$ for the apparatus, and compute
the Nexus coupling between $\bar\rho_{\rm app}$ and the
particle hidden variable $\lambda$.
\end{openprob}

\begin{openprob}[Amplitudes $A_5, A_3$ of $f_{H_4}$]
\label{op:amplitudes}
Derive the leading Fourier amplitudes $A_5$ (5-fold) and $A_3$
(3-fold) of $f_{H_4}(\theta)$ from the $H_4$ symmetry of the
600-cell adjacency matrix.  These amplitudes set the relative
magnitudes of the two angular components of the Bell-test deviation,
and are the key discriminating prediction between CPP and other
hidden-variable theories.
\end{openprob}

\begin{openprob}[Decoherence threshold for Nexus visibility]
\label{op:decoherence}
At what decoherence time $\tau_{\rm dec}$ does the Nexus
correlation $K$ become detectable?  The CPP estimate $\eps_{\rm QP}
\sim N_q / N_{\rm sub}$ suggests that a quantum processor with
$N_q \to N_{\rm sub}$ qubits would show $\eps \to 1$ (maximal Nexus
effects), but decoherence would eliminate the quantum state before
this regime is reached.  Find the $N_q$--$\tau_{\rm dec}$ trade-off
curve and identify the optimal quantum architecture for testing
CPP.
\end{openprob}

%=====================================================================
\section{Position in the CPP Programme}
%=====================================================================

This companion paper (SD\#1) is explicitly a \emph{position paper and
open problem statement}, not a solved derivation.  It occupies the
same role in the quantum foundations programme that companion C14
(confinement from qDP chaining \citep{c14}) played before the strong
sector series: it identifies the mechanism, states the open problem,
and provides the predictive framework that the subsequent derivation
must deliver.

The priority order for future CPP work is:

\begin{center}
\begin{tabular}{@{}lll@{}}
\toprule
Priority & Task & Why \\
\midrule
1 & Open Problem~\ref{op:K}: derive $K$ & Core superdeterminism result \\
2 & Open Problem~\ref{op:apparatus}: apparatus model & Required input for $K$ \\
3 & Open Problem~\ref{op:amplitudes}: $A_5, A_3$ & Key discriminating prediction \\
4 & Open Problem~\ref{op:decoherence}: decoherence threshold & Experimental target \\
5 & EW Open Problem EW-1: derive $\eta$ & Closes EW mass sector \\
6 & SS Open Problem SS-2: derive $\sigma$ & Closes strong mass sector \\
\bottomrule
\end{tabular}
\end{center}

The derivation of $K$ (priority 1) is the highest-value outstanding
problem in the CPP programme because it is the only one that could
produce a \emph{novel, falsifiable prediction} distinguishable from
standard QM.  All other open problems close existing derivations;
this one opens a new experimental window.

%=====================================================================
\section{Conclusion}
%=====================================================================

The CPP programme instantiates the superdeterministic hidden-variable
framework that Hossenfelder \citep{hossenfelder2024} argues is the
most promising path beyond standard quantum mechanics.  The DP~Sea
is the hidden-variable space; the Nexus is the superdeterministic
mechanism; the 600-cell is the geometric structure that fixes the
symmetry of the deviations from standard QM.

The central open problem is the derivation of the Nexus correlation
function $K(\lambda, \theta_A, \theta_B)$.  Three constraints on $K$
are already established: it has zero mean over $\lambda$, its
magnitude is bounded by the participation fraction
$\eps \sim N_{\rm part}/N_{\rm app}$, and its angular dependence is
governed by the $H_4$ symmetry of the 600-cell.

The primary experimental prediction is the \emph{600-cell angular
signature}: a residual in precision Bell-test scans of $E(\theta)$
with 5-fold and 3-fold harmonic components, with extrema at the
golden-ratio angles $\{31.7°, 36°, 58.3°, 72°, 120°\}$ of the
icosahedral subgroup of $H_4$.  The amplitude is $\sim 10^{-26}$
for standard Bell tests and scales with $N_{\rm part}/N_{\rm app}$,
potentially reaching $\sim 10^{-20}$ in megaqubit quantum processors.

This prediction is specific, non-trivial, and falsifiable.  If
quantum computers produce the anomalies that Hossenfelder predicts,
CPP specifies exactly what pattern those anomalies should show.

\bibliographystyle{unsrtnat}
\begin{thebibliography}{14}

\bibitem{hossenfelder2024}
S.~Hossenfelder,
``Why faster than light travel is possible after all (and what this
has to do with quantum physics),''
YouTube lecture transcript (2024).

\bibitem{cpp_qm_series}
T.~L. Abshier and Grok,
``Quantum mechanics from the 600-cell lattice,''
CPP QM Series 2040a--f v3.1, Hyperphysics Institute (2026).

\bibitem{cpp_ew_unified}
T.~L. Abshier and Grok,
``The electroweak sector from the 600-cell lattice,''
CPP EW Unified v2.1, Hyperphysics Institute (2026).

\bibitem{cpp_ss_unified}
T.~L. Abshier and Grok,
``The strong sector from the 600-cell lattice,''
CPP SS Unified v2, Hyperphysics Institute (2026).

\bibitem{cpp2040c}
T.~L. Abshier and Grok,
``Born rule and measurement from the DP Sea,''
CPP-2040c v3.1, Hyperphysics Institute (2026).

\bibitem{cpp2040d}
T.~L. Abshier and Grok,
``Open quantum systems and the Lindblad equation,''
CPP-2040d v3.1, Hyperphysics Institute (2026).

\bibitem{c14}
T.~L. Abshier and Grok,
``Quark confinement and the Cornell potential from qDP chaining,''
CPP companion~C14, Hyperphysics Institute (2026).

\bibitem{pdg2026}
Particle Data Group,
``Review of Particle Physics 2026,''
\textit{Phys.\ Rev.\ D} \textbf{110}, 030001 (2026).

\bibitem{bell1964}
J.~S. Bell,
``On the Einstein--Podolsky--Rosen paradox,''
\textit{Physics} \textbf{1}, 195 (1964).

\bibitem{hossenfelder_palmer}
S.~Hossenfelder and T.~Palmer,
``Rethinking superdeterminism,''
\textit{Front.\ Phys.} \textbf{8}, 139 (2020).

\bibitem{li2018}
M.-H. Li \textit{et al.},
``Test of local realism into the past without detection and
locality loopholes,''
\textit{Phys.\ Rev.\ Lett.} \textbf{121}, 080404 (2018).

\bibitem{hensen2015}
B.~Hensen \textit{et al.},
``Loophole-free Bell inequality violation using electron spins,''
\textit{Nature} \textbf{526}, 682 (2015).

\bibitem{brouwer1989}
A.~E. Brouwer, A.~M. Cohen, and A.~Neumaier,
\textit{Distance-Regular Graphs}, Springer (1989).

\bibitem{abshier_sr}
T.~L. Abshier,
``Mechanistic derivation of relativistic effects via Space Stress
Vector,''
viXra:2511.0062 (2025).

\end{thebibliography}

\end{document}
